% GENERATED FILE. Edit canonical lessons, then run npm run generate:books.
% canonical-source-hash: fnv1a64:d294b519b09f077d
% canonical-lessons: RU-C480-khokhotat, RU-C480-khikhikat, RU-C480-rydat, RU-C480-uteshat, RU-C480-utesheniye

\chapter{Laughter and Tears}
\label{ch:ru-laughter-and-tears}
\hlchaptermodality{480}

\begin{chapteropening}
\textbf{By the end of this chapter:} \emph{I can describe ways of laughing and crying.}

Complete the last lesson of chapter 480: I can describe ways of laughing and crying.
\end{chapteropening}

\section[\texorpdfstring{\ru{хохотать} (khokhotát)}{хохотать (khokhotát)}]{\ru{хохотать} (khokhotát) — to roar with laughter}

\label{lesson:RU-C480-khokhotat}



\begin{quote}
Before the new one: say the Russian for to respect, then the Russian for sincere.
\end{quote}

\subsection*{You'll want to know: \ru{хохотать}}
\textbf{\ru{хохотать}} (\emph{khokhotát}) — \textquotedblleft{}to roar with laughter\textquotedblright{}.

\begin{quote}
\ru{Мы} \ru{хохотали} \ru{весь} \ru{вечер}.
\end{quote}

\begin{quote}
\emph{My khokhotáli ves vécher.}
\end{quote}

\begin{quote}
We roared with laughter all evening.
\end{quote}

Imperfective; present \ru{хохочу}, \ru{хохочешь} (\ru{т} becomes \ru{ч}). Perfective \ru{захохотать} — to burst out laughing.

\begin{cousinweb}
\emph{\ru{хохотать}} is inherited from Proto-Slavic \emph{xoxotati} (EtymDB), an imitative word built on the sound of loud laughter.
\end{cousinweb}

\begin{grammarlens}[title={What you've built}]
One more word for telling, explaining and describing at B1.
\end{grammarlens}

\subsection*{Your turn}
\begin{itemize}
\raggedright
  \item \emph{Say it:} \emph{khokhotát}
  \item \emph{Say it:} \emph{khokhotát}, once more
  \item \emph{Say it:} a sentence of your own with \emph{khokhotát}
  \item \emph{Recall:} say the Russian for to respect, then the Russian for sincere, then say \emph{khokhotát} again
\end{itemize}

\begin{tcolorbox}[breakable,colback=teal!4,colframe=teal!35!black,title={Before you move on}]
What does \ru{хохотать} mean? (\textquotedblleft{}To roar with laughter\textquotedblright{}.) Say it once more.
\end{tcolorbox}
\section[\texorpdfstring{\ru{хихикать} (khikhíkat)}{хихикать (khikhíkat)}]{\ru{хихикать} (khikhíkat) — to giggle}

\label{lesson:RU-C480-khikhikat}



\begin{quote}
Before the new one: say the Russian for sincere, then the Russian for to roar with laughter.
\end{quote}

\subsection*{You'll want to know: \ru{хихикать}}
\textbf{\ru{хихикать}} (\emph{khikhíkat}) — \textquotedblleft{}to giggle\textquotedblright{}.

\begin{quote}
\ru{Девочки} \ru{хихикают} \ru{на} \ru{последней} \ru{парте}.
\end{quote}

\begin{quote}
\emph{Dévochki khikhíkayut na poslédney párte.}
\end{quote}

\begin{quote}
The girls are giggling at the back desk.
\end{quote}

Imperfective; perfective \ru{хихикнуть} (once).

\begin{cousinweb}
\emph{\ru{хихикать}} is built on the imitative sound \emph{\ru{хи}-\ru{хи}}; EtymDB lists the word without a source.
\end{cousinweb}

\begin{grammarlens}[title={What you've built}]
One more word for telling, explaining and describing at B1.
\end{grammarlens}

\subsection*{Your turn}
\begin{itemize}
\raggedright
  \item \emph{Say it:} \emph{khikhíkat}
  \item \emph{Say it:} \emph{khikhíkat}, once more
  \item \emph{Say it:} a sentence of your own with \emph{khikhíkat}
  \item \emph{Recall:} say the Russian for sincere, then the Russian for to roar with laughter, then say \emph{khikhíkat} again
\end{itemize}

\begin{tcolorbox}[breakable,colback=teal!4,colframe=teal!35!black,title={Before you move on}]
What does \ru{хихикать} mean? (\textquotedblleft{}To giggle\textquotedblright{}.) Say it once more.
\end{tcolorbox}
\section[\texorpdfstring{\ru{рыдать} (rydát)}{рыдать (rydát)}]{\ru{рыдать} (rydát) — to sob}

\label{lesson:RU-C480-rydat}



\begin{quote}
Before the new one: say the Russian for to roar with laughter, then the Russian for to giggle.
\end{quote}

\subsection*{You'll want to know: \ru{рыдать}}
\textbf{\ru{рыдать}} (\emph{rydát}) — \textquotedblleft{}to sob\textquotedblright{}.

\begin{quote}
\ru{После} \ru{фильма} \ru{она} \ru{рыдала}.
\end{quote}

\begin{quote}
\emph{Pósle fílma oná rydála.}
\end{quote}

\begin{quote}
After the film she sobbed.
\end{quote}

Imperfective; present \ru{рыдаю}. Stronger than \emph{\ru{плакать}}.

\begin{cousinweb}
EtymDB lists \emph{\ru{рыдать}} without a source; it is usually traced to Proto-Slavic \emph{rydati} 'to weep aloud'.
\end{cousinweb}

\begin{grammarlens}[title={What you've built}]
One more word for telling, explaining and describing at B1.
\end{grammarlens}

\subsection*{Your turn}
\begin{itemize}
\raggedright
  \item \emph{Say it:} \emph{rydát}
  \item \emph{Say it:} \emph{rydát}, once more
  \item \emph{Say it:} a sentence of your own with \emph{rydát}
  \item \emph{Recall:} say the Russian for to roar with laughter, then the Russian for to giggle, then say \emph{rydát} again
\end{itemize}

\begin{tcolorbox}[breakable,colback=teal!4,colframe=teal!35!black,title={Before you move on}]
What does \ru{рыдать} mean? (\textquotedblleft{}To sob\textquotedblright{}.) Say it once more.
\end{tcolorbox}
\section[\texorpdfstring{\ru{утешать} (uteshát)}{утешать (uteshát)}]{\ru{утешать} (uteshát) — to comfort}

\label{lesson:RU-C480-uteshat}



\begin{quote}
Before the new one: say the Russian for to giggle, then the Russian for to sob.
\end{quote}

\subsection*{You'll want to know: \ru{утешать}}
\textbf{\ru{утешать}} (\emph{uteshát}) — \textquotedblleft{}to comfort\textquotedblright{}.

\begin{quote}
\ru{Мама} \ru{утешает} \ru{плачущего} \ru{ребёнка}.
\end{quote}

\begin{quote}
\emph{Máma utesháyet pláchushchego rebyónka.}
\end{quote}

\begin{quote}
Mum is comforting the crying child.
\end{quote}

Imperfective; perfective \ru{утешить}.

\begin{cousinweb}
\emph{\ru{утешать}} is built on \emph{\ru{у}-} and \emph{\ru{тешить}} 'to amuse, to please'; EtymDB lists \emph{\ru{тешить}} without a source, and it is usually traced to Proto-Slavic \emph{těšiti}.
\end{cousinweb}

\begin{grammarlens}[title={What you've built}]
One more word for telling, explaining and describing at B1.
\end{grammarlens}

\subsection*{Your turn}
\begin{itemize}
\raggedright
  \item \emph{Say it:} \emph{uteshát}
  \item \emph{Say it:} \emph{uteshát}, once more
  \item \emph{Say it:} a sentence of your own with \emph{uteshát}
  \item \emph{Recall:} say the Russian for to giggle, then the Russian for to sob, then say \emph{uteshát} again
\end{itemize}

\begin{tcolorbox}[breakable,colback=teal!4,colframe=teal!35!black,title={Before you move on}]
What does \ru{утешать} mean? (\textquotedblleft{}To comfort\textquotedblright{}.) Say it once more.
\end{tcolorbox}
\section[\texorpdfstring{\ru{утешение} (uteshéniye)}{утешение (uteshéniye)}]{\ru{утешение} (uteshéniye) — consolation, comfort}

\label{lesson:RU-C480-utesheniye}



\begin{quote}
Before the new one: say the Russian for to sob, then the Russian for to comfort.
\end{quote}

\subsection*{You'll want to know: \ru{утешение}}
\textbf{\ru{утешение}} (\emph{uteshéniye}) — \textquotedblleft{}consolation, comfort\textquotedblright{}.

\begin{quote}
\ru{Твои} \ru{слова} — \ru{для} \ru{меня} \ru{большое} \ru{утешение}.
\end{quote}

\begin{quote}
\emph{Tvoí slová — dlya menyá bolshóye uteshéniye.}
\end{quote}

\begin{quote}
Your words are a great comfort to me.
\end{quote}

Neuter.

\begin{cousinweb}
\emph{\ru{утешение}} is built on \emph{\ru{утешить}} 'to comfort' with \emph{-\ru{ение}}; the root is that of \emph{\ru{тешить}} 'to please', usually traced to Proto-Slavic \emph{těšiti}.
\end{cousinweb}

\begin{grammarlens}[title={What you've built}]
One more word for telling, explaining and describing at B1.
\end{grammarlens}

\subsection*{Your turn}
\begin{itemize}
\raggedright
  \item \emph{Say it:} \emph{uteshéniye}
  \item \emph{Say it:} \emph{uteshéniye}, once more
  \item \emph{Say it:} a sentence of your own with \emph{uteshéniye}
  \item \emph{Recall:} say the Russian for to sob, then the Russian for to comfort, then say \emph{uteshéniye} again
\end{itemize}

\begin{tcolorbox}[breakable,colback=teal!4,colframe=teal!35!black,title={Before you move on}]
What does \ru{утешение} mean? (\textquotedblleft{}Consolation, comfort\textquotedblright{}.) Say it once more.
\end{tcolorbox}
